\documentclass{article}
\usepackage[T2A]{fontenc}        
\usepackage[utf8]{inputenc}    
\usepackage[russian,english]{babel}
\usepackage{graphicx}
\usepackage{float}
\usepackage[hidelinks]{hyperref}
\usepackage{fancyvrb}

\title{Computer Skills for Scientific Writing\\Работа с графикой в LaTeX}
\author{Абрамян Артём Арменович}
\date{Октябрь 2025}

\begin{document}
  \maketitle
  \pagebreak

  \textbf{Цель работы:} Изучить возможности включения и манипулирования графикой в LaTeX, освоить различные способы позиционирования изображений, научиться работать с перекрестными ссылками.\smallskip

  \section{Введение}
  \label{sec:intro}
  
  В данной работе мы изучаем работу с графикой в LaTeX согласно документации. Для включения графики используется пакет \texttt{graphicx}, который добавляет команду \texttt{\textbackslash includegraphics}.

  \section{Базовое включение графики}
  \label{sec:basic}
  
  Создаем простой LaTeX файл с графикой. Пакет \texttt{graphicx} поддерживает форматы EPS, PNG, JPG и PDF.

  \begin{figure}[!h]
    \centering
    
\includegraphics[width=0.8\textwidth]{fig01.png}
    \caption{Базовый пример включения изображения}
    \label{fig:fig01}
  \end{figure}

  На рисунке~\ref{fig:fig01} показан базовый код для включения изображения в документ.

  \section{Манипулирование размерами графики}
  \label{sec:dimensions}

  \subsection{Использование параметра width}
  \label{subsec:width}
  
  Параметр \texttt{width} позволяет задать ширину изображения. Можно использовать абсолютные значения (например, 5cm) или относительные (например, 0.5\textbackslash textwidth).

  Пример кода:

  \texttt{\textbackslash includegraphics[width=0.5\textbackslash textwidth]\{image.png\}}

  \begin{figure}[!h]
    \centering
    \includegraphics[width=0.8\textwidth]{fig02.png}
    \caption{Изображение с параметром width=0.5\textbackslash textwidth}
    \label{fig:fig02}
  \end{figure}

  \subsection{Использование параметра height}
  \label{subsec:height}
  
  Аналогично ширине, можно задать высоту изображения с помощью параметра \texttt{height}. LaTeX автоматически сохраняет пропорции изображения.

  Пример кода:

  \texttt{\textbackslash includegraphics[height=5cm]\{image.png\}}

  \begin{figure}[!h]
    \centering
    \includegraphics[height=5cm]{fig03.png}
    \caption{Изображение с параметром height=5cm}
    \label{fig:fig03}
  \end{figure}

  \subsection{Параметр scale}
  \label{subsec:scale}
  
  Параметр \texttt{scale} позволяет масштабировать изображение пропорционально. Например, scale=0.5 уменьшит изображение в два раза.

  Пример кода:

  \texttt{\textbackslash includegraphics[scale=0.5]\{image.png\}}

  \begin{figure}[!h]
    \centering
    \includegraphics[scale=0.5]{fig04.png}
    \caption{Использование параметра scale=0.5}
    \label{fig:fig04}
  \end{figure}

  \subsection{Поворот изображения}
  \label{subsec:angle}
  
  С помощью параметра \texttt{angle} можно повернуть изображение на заданный угол в градусах.

  Пример кода:

  \texttt{\textbackslash includegraphics[angle=45,width=0.6\textbackslash textwidth]\{image.png\}}

  \begin{figure}[!h]
    \centering
    \includegraphics[angle=45,width=0.6\textwidth]{fig05.png}
    \caption{Поворот изображения на 45 градусов с помощью параметра angle=45}
    \label{fig:fig05}
  \end{figure}

  \section{Разница между \textbackslash textwidth и \textbackslash linewidth}
  \label{sec:difference}
  
  \texttt{\textbackslash textwidth} --- это ширина текстового блока на физической странице, тогда как \texttt{\textbackslash linewidth} --- это текущая ширина, которая может локально отличаться.

  \begin{figure}[!h]
    \centering
    \includegraphics[width=0.8\textwidth]{fig06.png}
    \caption{Пример кода с использованием \textbackslash textwidth}
    \label{fig:fig06}
  \end{figure}

  Пример использования \texttt{\textbackslash textwidth}:

  \texttt{\textbackslash includegraphics[width=0.5\textbackslash textwidth]\{image.png\}}

  Пример использования \texttt{\textbackslash linewidth}:

  \texttt{\textbackslash includegraphics[width=0.5\textbackslash linewidth]\{image.png\}}

  Разница наиболее заметна при использовании опции класса \texttt{twocolumn}. В режиме двух колонок \texttt{\textbackslash linewidth} адаптируется к ширине колонки, а \texttt{\textbackslash textwidth} всегда относится к полной ширине страницы.

  \begin{figure}[!h]
    \centering
    \includegraphics[width=0.8\textwidth]{fig07.png}
    \caption{Демонстрация разницы в режиме twocolumn}
    \label{fig:fig07}
  \end{figure}

  В обычном однокалоночном режиме эти значения одинаковы, но в \texttt{twocolumn} изображение с \texttt{width=0.5\textbackslash linewidth} будет занимать половину ширины колонки, а с \texttt{width=0.5\textbackslash textwidth} --- половину всей страницы.

  \section{Позиционирование флоатов}
  \label{sec:floats}
  
  В LaTeX изображения обычно включаются как флоаты (floats), чтобы избежать больших пустых промежутков на странице.

  \subsection{Спецификаторы позиции}
  \label{subsec:specifiers}
  
  Можно использовать до четырех спецификаторов позиции:
  \begin{itemize}
    \item \texttt{h} --- 'Here' (здесь, если возможно)
    \item \texttt{t} --- Top (верх страницы)
    \item \texttt{b} --- Bottom (низ страницы)
    \item \texttt{p} --- Page (отдельная страница только для флоатов)
  \end{itemize}

  Пример использования спецификатора [ht]:

  \texttt{\textbackslash begin\{figure\}[ht]}

  \texttt{\textbackslash centering}

  \texttt{\textbackslash includegraphics[width=0.5\textbackslash textwidth]\{image.png\}}

  \texttt{\textbackslash caption\{Изображение\}}

  \texttt{\textbackslash end\{figure\}}

  \begin{figure}[!h]
    \centering
    \includegraphics[width=0.8\textwidth]{fig08.png}
    \caption{Пример кода с различными спецификаторами позиции}
    \label{fig:fig08}
  \end{figure}

  \subsection{Использование пакета float}
  \label{subsec:float_package}
  
  Пакет \texttt{float} добавляет спецификатор \texttt{H}, который размещает изображение точно в месте его определения в коде.

  Сначала подключаем пакет:

  \texttt{\textbackslash usepackage\{float\}}

  Затем используем спецификатор [H]:

  \texttt{\textbackslash begin\{figure\}[H]}

  \texttt{\textbackslash centering}

  \texttt{\textbackslash includegraphics[width=0.5\textbackslash textwidth]\{image.png\}}

  \texttt{\textbackslash caption\{Изображение точно здесь\}}

  \texttt{\textbackslash end\{figure\}}

  \begin{figure}[!h]
    \centering
    \includegraphics[width=0.8\textwidth]{fig09.png}
    \caption{Код с использованием спецификатора [H] из пакета float}
    \label{fig:fig09}
  \end{figure}

%   \begin{figure}[!h]
%     \centering
%     \includegraphics[width=0.8\textwidth]{fig10.png}
%     \caption{Результат компиляции с флоатами}
%     \label{fig:fig10}
%   \end{figure}

  \section{Перекрестные ссылки}
  \label{sec:crossref}
  
  \subsection{Механизм \textbackslash label и \textbackslash ref}
  \label{subsec:label_ref}
  
  Для автоматической нумерации элементов используется механизм меток. Чтобы LaTeX запомнил место в документе, его нужно пометить с помощью \texttt{\textbackslash label}, а затем сослаться на него с помощью \texttt{\textbackslash ref}.

  Пример использования для разделов:

  \texttt{\textbackslash section\{Введение\}}

  \texttt{\textbackslash label\{sec:intro\}}

  \texttt{Как описано в разделе\textasciitilde\textbackslash ref\{sec:intro\}...}

  Пример использования для рисунков:

  \texttt{\textbackslash begin\{figure\}[ht]}

  \texttt{\textbackslash centering}

  \texttt{\textbackslash includegraphics[width=0.5\textbackslash textwidth]\{image.png\}}

  \texttt{\textbackslash caption\{Пример\}}

  \texttt{\textbackslash label\{fig:example\}}

  \texttt{\textbackslash end\{figure\}}

  \texttt{На рисунке\textasciitilde\textbackslash ref\{fig:example\} показано...}

  \begin{figure}[!h]
    \centering
    \includegraphics[width=0.8\textwidth]{fig11.png}
    \caption{Пример использования \textbackslash label и \textbackslash ref}
    \label{fig:fig11}
  \end{figure}

  Например, это подраздел~\ref{subsec:label_ref}, который находится в разделе~\ref{sec:crossref}. Ссылка на рисунок~\ref{fig:fig11} автоматически подставляет правильный номер.

  \subsection{Правильное размещение \textbackslash label}
  \label{subsec:correct_label}
  
  Команда \texttt{\textbackslash label} всегда ссылается на предыдущую нумерованную сущность. Это означает, что для флоатов \texttt{\textbackslash label} должна идти \textbf{после} команды \texttt{\textbackslash caption}.

  Правильный пример:

  \texttt{\textbackslash begin\{figure\}[ht]}

  \texttt{\textbackslash centering}

  \texttt{\textbackslash includegraphics[width=0.5\textbackslash textwidth]\{image.png\}}

  \texttt{\textbackslash caption\{Описание изображения\}}

  \texttt{\textbackslash label\{fig:correct\}}

  \texttt{\textbackslash end\{figure\}}

  \begin{figure}[!h]
    \centering
    \includegraphics[width=0.8\textwidth]{fig12.png}
    \caption{Правильное размещение метки --- после caption}
    \label{fig:fig12}
  \end{figure}

  \subsection{Неправильное размещение метки}
  \label{subsec:incorrect_label}
  
  Если поместить \texttt{\textbackslash label} перед \texttt{\textbackslash caption}, метка будет ссылаться на номер секции, а не на номер рисунка.

  Неправильный пример:

  \texttt{\textbackslash begin\{figure\}[ht]}

  \texttt{\textbackslash centering}

  \texttt{\textbackslash label\{fig:incorrect\}}

  \texttt{\textbackslash includegraphics[width=0.5\textbackslash textwidth]\{image.png\}}

  \texttt{\textbackslash caption\{Описание изображения\}}

  \texttt{\textbackslash end\{figure\}}

  \begin{figure}[!h]
    \centering
    \includegraphics[width=0.8\textwidth]{fig13.png}
    \caption{Демонстрация неправильного размещения метки}
    \label{fig:fig13}
  \end{figure}


  \subsection{Метки для уравнений}
  \label{subsec:equations}
  
  Для уравнений метка должна находиться внутри окружения \texttt{equation}, до команды \texttt{\textbackslash end\{equation\}}.

  \begin{equation}
  E = mc^2
  \label{eq:correct}
  \end{equation}

  Правильный пример:

  \texttt{\textbackslash begin\{equation\}}

  \texttt{E = mc\textasciicircum 2}

  \texttt{\textbackslash label\{eq:einstein\}}

  \texttt{\textbackslash end\{equation\}}

  \texttt{Уравнение\textasciitilde\textbackslash ref\{eq:einstein\} показывает...}

  Неправильный пример:

  \texttt{\textbackslash begin\{equation\}}

  \texttt{a\textasciicircum 2 + b\textasciicircum 2 = c\textasciicircum 2}

  \texttt{\textbackslash end\{equation\}}

  \texttt{\textbackslash label\{eq:wrong\}}

  \texttt{Уравнение\textasciitilde\textbackslash ref\{eq:wrong\} будет ссылаться на номер раздела!}

  \begin{figure}[!h]
    \centering
    \includegraphics[width=0.8\textwidth]{fig14.png}
    \caption{Правильное и неправильное размещение меток для уравнений}
    \label{fig:fig14}
  \end{figure}

  Если метка размещена после \texttt{\textbackslash end\{equation\}}, она будет ссылаться на номер раздела, а не на номер уравнения.

  \subsection{Количество компиляций}
  \label{subsec:compilation}
  
  Из-за использования вспомогательного файла, первая компиляция документа может показать вопросительные знаки (??) вместо номеров ссылок.

  \begin{figure}[!h]
    \centering
    \includegraphics[width=0.8\textwidth]{fig15.png}
    \caption{Первая компиляция --- вопросительные знаки вместо номеров}
    \label{fig:fig15}
  \end{figure}

  \begin{figure}[!h]
    \centering
    \includegraphics[width=0.8\textwidth]{fig16.png}
    \caption{Вторая компиляция --- правильные номера ссылок}
    \label{fig:fig16}
  \end{figure}

  После второй компиляции все ссылки будут правильными, как показано на рисунках~\ref{fig:fig15} и~\ref{fig:fig16}.

  \section{Использование пакета hyperref}
  \label{sec:hyperref}
  
  Пакет \texttt{hyperref} позволяет превратить перекрестные ссылки в гиперссылки. В данном документе используется опция \texttt{hidelinks}, чтобы ссылки имели тот же цвет, что и обычный текст.

  Пример подключения:

  \texttt{\textbackslash usepackage[hidelinks]\{hyperref\}}

  Другие полезные опции:

  \texttt{\textbackslash usepackage[colorlinks=true,}

  \texttt{linkcolor=blue,}

  \texttt{citecolor=red]\{hyperref\}}

%   \begin{figure}[!h]
%     \centering
%     \includegraphics[width=0.8\textwidth]{fig17.png}
%     \caption{Подключение пакета hyperref}
%     \label{fig:fig17}
%   \end{figure}

  Теперь можно кликнуть на любую ссылку, например на раздел~\ref{sec:intro}, и перейти к соответствующему месту в документе.

  \section{Заключение}
  \label{sec:conclusion}
  
  В данной работе были продемонстрированы:
  \begin{itemize}
    \item Базовое включение графики в документ (раздел~\ref{sec:basic})
    \item Манипулирование размерами изображений с помощью width, height, scale и angle (раздел~\ref{sec:dimensions})
    \item Разница между \textbackslash textwidth и \textbackslash linewidth (раздел~\ref{sec:difference})
    \item Различные спецификаторы позиционирования флоатов (раздел~\ref{sec:floats})
    \item Правильное использование механизма меток и ссылок (раздел~\ref{sec:crossref})
    \item Создание гиперссылок с помощью пакета hyperref (раздел~\ref{sec:hyperref})
  \end{itemize}

  \textbf{Выводы:} В данной работе мы изучили продвинутые возможности работы с графикой в LaTeX. Освоили управление размерами и позиционированием изображений, а также правильное использование системы перекрестных ссылок. Особое внимание было уделено распространенным ошибкам при размещении меток и способам их избежания. Выяснили, что для корректной работы перекрестных ссылок необходимо компилировать документ минимум два раза.

\end{document}